\documentclass[12pt,a4paper]{article}
        \makeatletter
        \def\input@path{{../}}
        \makeatother
        \usepackage{almeria-fichas}

        \FichaMeta{Sesion 07}{Perimetro de zonas urbanas}{Bloque 2 · Geometria urbana}{Calcular el perimetro de barrios, plazas o parques modelizados como poligonos simples.}{Procedimientos}{Comprender, Aplicar, Analizar}

        \begin{document}

        \FichaCabecera

        \begin{materialesbox}
        \begin{itemize}
    \item Ficha impresa
    \item Lapiz
    \item Regla
    \item Calculadora sencilla si se necesita
\end{itemize}
        \end{materialesbox}

        \begin{pasosbox}
        \begin{enumerate}
    \item Rodea el contorno completo antes de sumar.
    \item Anota cada lado una sola vez.
    \item Escribe la operacion antes del resultado final.
    \item Comprueba la unidad al terminar.
\end{enumerate}
        \end{pasosbox}

        \begin{recordatoriobox}
        \begin{itemize}
    \item El perimetro es la longitud del borde.
    \item Se suma lado con lado.
    \item La unidad del perimetro suele ser m o km.
\end{itemize}
        \end{recordatoriobox}

        \begin{apoyobox}
        \begin{itemize}
    \item Puedes numerar los lados para no olvidarte de ninguno.
    \item Si una figura es regular, piensa si puedes multiplicar.
    \item No pongas m²: el perimetro no mide superficie.
\end{itemize}
        \end{apoyobox}

        \BloomSection{comprender}

\BloomExercise{comprender}{1}{Perimetro o no}{Explica por que para hallar el perimetro de un rectangulo se suman sus lados y no su interior.}{6}

\BloomExercise{comprender}{2}{Atajo matematico}{Explica cuando es util hacer $2\,(l+a)$ en lugar de sumar los cuatro lados uno a uno.}{5}

\BloomSection{aplicar}

\BloomExercise{aplicar}{1}{Calculo directo}{Calcula el perimetro de un rectangulo de $18\,m$ y $12\,m$.}{5}

\BloomExercise{aplicar}{2}{Parque poligonal}{Un parque tiene lados de $9\,m$, $14\,m$, $11\,m$ y $14\,m$. Calcula su perimetro y escribe la operacion completa.}{6}

\BloomSection{analizar}

\BloomExercise{analizar}{1}{Partes y relaciones}{Analiza un ejemplo relacionado con \emph{Perimetro de zonas urbanas}. Separa sus partes, indica cuales son relevantes y explica como se relacionan entre si.}{8}

\BloomExercise{analizar}{2}{Detecta errores o diferencias}{Compara dos respuestas, dos representaciones o dos procedimientos relacionados con esta sesion. Analiza sus diferencias y explica que errores o mejoras detectas.}{8}

        \begin{evidenciabox}
        \begin{itemize}
    \item Calcular perimetros con suma y con multiplicacion cuando conviene.
    \item Interpretar el perimetro como borde exterior.
    \item Explicar una estrategia de calculo clara.
\end{itemize}
        \end{evidenciabox}

        \end{document}
